\section{An elementary estimate for the four-vertex pattern}
\label{num:elementary-estimates}

The endpoint flux and projection identities reduce the four-vertex
configuration to the following inequality. Its variables are scalars;
no geometric assumptions beyond their displayed bounds are needed.

\begin{lemma}\label{num:four-vertex-scalar}
Suppose
\[
 \begin{gathered}
 0<d\le\frac1{\sqrt3},\qquad
 0\le x\le\sqrt\pi,\qquad 0\le\tau<1,\\
 \tau+\frac{x}{\pi}(1-\tau^2)=1-d^2.
 \end{gathered}
\]
Then
\[
 d+\frac{\tau^2+(\tau+xd)^2}{\pi}<1.
\]
\end{lemma}

\begin{proof}
Write $R(\tau,d)$ for the expression to be bounded, with
$x=\pi(1-\tau-d^2)/(1-\tau^2)$. In particular
$\tau\le1-d^2$. Let $f(q)=q+x(1-q^2)/\pi$.
Suppose for contradiction that $R(\tau,d)\ge1$.
All terminating decimals below denote exact rational numbers; the
elementary bounds $157/50<\pi<22/7$ suffice.

First, $\tau>1/10$.  Otherwise the increasing function
$q+(1-q^2)/\sqrt\pi$ on $[0,1/10]$ would give
\[
 \frac23\le f(\tau)
 \le\frac1{10}+\frac{99}{100\sqrt\pi}<\frac23.
\]
For the last inequality, square the positive terms and use
$(99/100)^2/(157/50)<(17/30)^2$.

\smallskip
\noindent\emph{The range $1/10<\tau\le3/10$.}
Hold $\tau$ fixed and increase $d$ to $1/\sqrt3$, redefining $x$ by
the formula above.  Along this comparison $x$ decreases and
remains nonnegative, because
$1-\tau-d^2\ge1-3/10-1/3>0$.  Thus $x\le\sqrt\pi$ remains valid.
Direct differentiation gives
\[
 \frac{\partial R}{\partial d}
 =1+\frac{2(\tau+xd)(1-\tau-3d^2)}{1-\tau^2}.
\]
If the numerator's last factor is nonnegative this derivative is at
least one.  Otherwise use $1-\tau-3d^2\ge-\tau$,
$1-\tau^2\ge91/100$, and
\[
 \tau+xd\le\frac3{10}+\sqrt{\frac\pi3}
 <\frac3{10}+\frac{41}{40}=\frac{53}{40}
\]
to obtain
\[
 \frac{\partial R}{\partial d}
 >1-\frac{2(3/10)(53/40)}{91/100}
 =\frac{23}{182}>0.
\]
It remains to bound $R(\tau,1/\sqrt3)$.  Set
\[
 J(\tau)=\tau+
 \frac{\pi(2/3-\tau)}{\sqrt3(1-\tau^2)}.
\]
On $0\le\tau\le3/10$,
\begin{align*}
 &1+\tau^2-\frac43\tau-\frac45(1-\tau^2)^2\\
 &\hspace{1cm}=
 \frac15-\frac43\tau+\frac{13}{5}\tau^2-\frac45\tau^4
 \ge\frac15-\frac43\tau+\frac{316}{125}\tau^2>0.
\end{align*}
The last quadratic has positive leading coefficient and discriminant
$-1376/5625$.  Hence
\[
 J'(\tau)
 =1-\frac\pi{\sqrt3}
     \frac{1+\tau^2-4\tau/3}{(1-\tau^2)^2}
 <1-\frac45\frac95=-\frac{11}{25},
\]
where $\pi/\sqrt3>9/5$ follows by squaring.  Also
\[
 J(3/10)=\frac3{10}+\frac{110\pi}{273\sqrt3}
 >\frac3{10}+\frac{66}{91}>1.
\]
Thus $J>1$ throughout this range and
$(J^2+\tau^2)'<-22/25+3/5<0$.
For an exact endpoint bound, use $1/\sqrt3<26/45$ to get
\[
 J(1/10)
 <\frac1{10}+\frac{22}{7}\frac{26}{45}\frac{170}{297}
 =\frac{19381}{17010}<\frac{57}{50}.
\]
It follows that
\[
 R(\tau,d)
 \le R(1/10,1/\sqrt3)
 <\frac{26}{45}
   +\frac{1/100+(57/50)^2}{157/50}
 =\frac{35143}{35325}<1.
\]

\smallskip
\noindent\emph{The range $3/10\le\tau<1$.}
Put $U=\sqrt{1-\tau}$ and $Z=d/U$.  The inequality
$d^2\le1-\tau$ gives $0\le Z\le1$, while
$0<U\le\sqrt{7/10}$.  The formula for $x$ becomes
$x=\pi(1-Z^2)/(1+\tau)$, and
the contradictory assumption $R(\tau,d)\ge1$ becomes
\begin{equation}\label{slice:four-rescaled}
 \frac{1-2\tau^2/\pi}{\sqrt{1-\tau}}
 \le Z+\frac{2\tau}{1+\tau}Z(1-Z^2)
       +\frac{\pi U}{(1+\tau)^2}Z^2(1-Z^2)^2.
\end{equation}
The derivative of the left side is
\[
 \frac{1/2+(3\tau^2-4\tau)/\pi}{(1-\tau)^{3/2}}>0,
\]
because $3\tau^2-4\tau\ge-4/3$ and $\pi>8/3$.
Using $\sqrt{7/10}<837/1000$, its value at $3/10$ exceeds
\[
 \frac{148000}{131409}>
 \frac98=1.125.
\]

To bound the right side, put $c_*=2\pi/(3\sqrt3)$.  We claim that
\begin{equation}\label{slice:four-coefficient}
 \frac{2\tau}{1+\tau}
 +\frac{2\pi U}{3\sqrt3(1+\tau)^2}<\frac{111}{100}.
\end{equation}
Multiplication by $(1+\tau)^2=(2-U^2)^2$ shows that this is equivalent
to positivity of
\[
 \frac{11}{25}-c_*U+\frac{39}{25}U^2-\frac{89}{100}U^4
 \ge\frac{11}{25}-c_*U+\frac{937}{1000}U^2.
\]
The last quadratic is positive for every real $U$: its leading
coefficient is positive and its discriminant satisfies
\[
 c_*^2-4\frac{11}{25}\frac{937}{1000}
 <\frac{1936}{1323}-\frac{10307}{6250}
 =-\frac{1536161}{8268750}<0.
\]
This proves \eqref{slice:four-coefficient}.  Since
$0\le Z(1-Z^2)\le2/(3\sqrt3)$, the right side of
\eqref{slice:four-rescaled} is at most
\[
 Z+\frac{111}{100}Z(1-Z^2).
\]
Its maximum on $[0,1]$ is
\[
 \frac{211}{150}\sqrt{\frac{211}{333}}<\frac98,
\]
as follows by squaring from
\[
 \frac{81}{64}-\frac{9393931}{7492500}
 =\frac{1420229}{119880000}>0.
\]
This contradicts \eqref{slice:four-rescaled} and completes both ranges.
The case $x=0$ was included: then
$\tau=1-d^2\in[2/3,1)$, in the second range.
\end{proof}
